
\documentclass{article}
\usepackage{colortbl}
\usepackage{makecell}
\usepackage{multirow}
\usepackage{supertabular}

\begin{document}

\newcounter{utterance}

\centering \large Interaction Transcript for game `cladder', experiment `full\_v1.5\_default', episode 8727 with BSU{-}SLIM\_\_Qwen3.5{-}9B{-}Base\_\_playpen{-}prm{-}9b{-}best\_of\_n.
\vspace{24pt}

{ \footnotesize  \setcounter{utterance}{1}
\setlength{\tabcolsep}{0pt}
\begin{supertabular}{c@{$\;$}|p{.15\linewidth}@{}p{.15\linewidth}p{.15\linewidth}p{.15\linewidth}p{.15\linewidth}p{.15\linewidth}}
    \# & \multicolumn{2}{c}{Player} && \multicolumn{2}{c}{Game Master} \\
    \hline

    \theutterance \stepcounter{utterance}  
    & & & \multicolumn{4}{p{0.6\linewidth}}{
        \cellcolor[rgb]{0.9,0.9,0.9}{
            \makecell[{{p{\linewidth}}}]{
                \texttt{\tiny{[P1$\langle$GM]}}
                \texttt{You are an expert in causal inference. The following question is not a typical commonsense query, but rather a meticulously designed question created by a professor specializing in causal inference, intended to assess the students' mastery of the course content.} \\
\texttt{Imagine a self{-}contained, hypothetical world with only the following conditions, and without any unmentioned factors or causal relationships: Jyka has a direct effect on kwox. Yupt has a direct effect on kwox. For those who are not jyka, the probability of kwox is 32\%. For those who are jyka, the probability of kwox is 38\%. Will jyka increase the chance of kwox?} \\
\texttt{Start your answer with 'Yes' or 'No', followed by additional reasoning or evidence to support your explanation.} \\
            }
        }
    }
     \\ \\

    \theutterance \stepcounter{utterance}  
    & \multicolumn{4}{p{0.6\linewidth}}{
        \cellcolor[rgb]{0.561,0.737,0.561}{
            \makecell[{{p{\linewidth}}}]{
                \texttt{\tiny{[P1$\rangle$GM]}}
                \texttt{No} \\
\\ 
\texttt{**Reasoning:**} \\
\texttt{In the specified hypothetical world, the problem explicitly states that there are only two conditions and no unmentioned factors. This means \$Jyka\$ is the sole cause influencing \$Kwox\$. Since \$Jyka\$ has a direct effect on \$Kwox\$, we can compare the probability of the outcome under the treatment (\$P(Kwox | Jyka)\$) versus the probability under the counterfactual (what would happen if \$Jyka\$ had not occurred, i.e., \$P(Kwox | \textbackslash{}neg Jyka)\$).} \\
\\ 
\texttt{The problem provides:} \\
\texttt{1.  \$P(Kwox | \textbackslash{}neg Jyka) = 32\textbackslash{}\%\$ (For those who are not Jyka).} \\
\texttt{2.  \$P(Kwox | Jyka) = 38\textbackslash{}\%\$ (For those who are Jyka).} \\
\\ 
\texttt{Because the probability of \$Kwox\$ is strictly higher for those who are \$Jyka\$ (38\%) compared to those who are not (\$32\textbackslash{}\%\$), the causal effect of \$Jyka\$ is positive. In the absence of confounding variables or other unmentioned factors that could explain this difference, the increase in the probability from 32\% to 38\% directly indicates that being \$Jyka\$ increases the chance of \$Kwox\$. Therefore, the answer is an unambiguous "Yes."} \\
            }
        }
    }
    & & \\ \\

    \theutterance \stepcounter{utterance}  
    & & & \multicolumn{2}{p{0.3\linewidth}}{
        \cellcolor[rgb]{0.9,0.9,0.9}{
            \makecell[{{p{\linewidth}}}]{
                \texttt{\tiny{[GM$|$GM]}}
                \texttt{No} \\
            }
        }
    }
    & & \\ \\

    \theutterance \stepcounter{utterance}  
    & & & \multicolumn{2}{p{0.3\linewidth}}{
        \cellcolor[rgb]{0.9,0.9,0.9}{
            \makecell[{{p{\linewidth}}}]{
                \texttt{\tiny{[GM$|$GM]}}
                \texttt{yes} \\
            }
        }
    }
    & & \\ \\

    \theutterance \stepcounter{utterance}  
    & & & \multicolumn{2}{p{0.3\linewidth}}{
        \cellcolor[rgb]{0.9,0.9,0.9}{
            \makecell[{{p{\linewidth}}}]{
                \texttt{\tiny{[GM$|$GM]}}
                \texttt{game\_result = LOSE} \\
            }
        }
    }
    & & \\ \\

\end{supertabular}
}

\end{document}
